\section{Outcomes, original-root estimands and inference}
\label{sec:b-statistics}

\subsection{Outcomes}

The primary outcome is $U$, the hidden-test accuracy of the final valid candidate
in $[0,1]$.  The following are kept distinct: $U$; the binary validated
completion $S$ and validity; individual degradation frequency; the measured
resource vector; execution status; and data-integrity status.

\begin{itemize}[leftmargin=1.4em]
\item \textbf{Invalid or no candidate at the horizon:} $U=0$ with the explicit
  invalid flag \texttt{outcome\_invalid}$=1$ and its \texttt{invalid\_reason}
  retained.
\item \textbf{Collection failure:} an integrity status
  (\texttt{ok}$\mid$\texttt{collection\_failure}$\mid$\texttt{infrastructure\_stop}),
  \emph{never} a utility value.  The evaluator refuses to analyse it as $0$.
\item \textbf{Unavailable action:} no observed outcome; the root follows its
  registered continue fallback and remains in the original-root analysis.
\end{itemize}

A contrast may not mix $S$ and $U$ as if they were one outcome.  $U$ is not
scored only on successful or valid branches, and a missing score caused by
collection failure is never interchangeable with task failure.  Utility
direction, normalization and clipping are frozen before test inspection.

\subsection{Original-root estimand}

For a policy $d$, the original-root value is
\begin{equation}
 \theta(d)=\E_{\text{original roots}}\big[\,U^{\,d,\text{fallback}}\,\big],
 \label{eq:original-root-value}
\end{equation}
where the expectation ranges over all original roots, including early-terminal
roots and roots at which the checkpoint action is infeasible, with the
registered fallback applied.  There is no conditioning on reaching checkpoint 2,
no transplanted checkpoint-2 state and at most one verification per root.  The
one-step and tree estimands are reported separately and never share a confidence
guarantee automatically.

\subsection{Primary contrasts and multiplicity}

The two primary contrasts are
\begin{equation}
 D_{\mathrm{myopic}}=\theta(d_{\mathrm{frozen}})-\theta(d_{\mathrm{myopic}}),
 \qquad
 D_{\mathrm{VOI}}=\theta(d_{\mathrm{frozen}})-\theta(d_{\mathrm{VOI}}),
 \label{eq:primary-contrasts}
\end{equation}
where $d_{\mathrm{VOI}}$ is the BAVAR-inspired context-dependent VOI scheduler
(\cref{sec:policy}), tested with \emph{simultaneous} intervals over the two
contrasts.  The sampling
unit is the template family; roots and sibling branches generated from one family
share their cluster and split.  A contrast is \emph{supported} when its adjusted
lower confidence bound exceeds zero; the practical target is $0.02\,U$; a
contrast whose adjusted upper bound is below $0.02$ rejects the practical target;
an interval spanning the target is \emph{inconclusive}, with no positive claim
and no automatic budget expansion.

\subsection{Allocation and budget conventions}

The primary capacity is $q=0.25$, with $q\in\{0.05,0.10,0.50\}$ descriptive
only.  The one-step lane preserves the $\lceil q\,n_{\text{feasible}}\rceil$
convention and reports effective capacity for small batches.  The original-root
tree evaluation registers $\lfloor q\,N_{\text{original}}\rfloor$ as the total
verifier-call budget with at most one call per root.  The two denominators are
never exchanged, and unspent capacity stays unspent.  Matched \emph{allowed}
budgets need not produce equal realized call counts; realized calls and work are
reported separately.

\subsection{Harm}

The mean no-injected-fault utility loss is at most $0.02$ versus never-verify;
this is an empirical endpoint, not a per-person no-harm guarantee and not a bound
on the probability of any degradation
(\texttt{safety\_claim\_level: empirical}).  Individual degradation frequency
and the invalid-output rate are reported separately.  An integrity violation
stops collection and is not covered by the utility harm allowance.

\subsection{Resampling and batch coupling}

Inference uses a template-family cluster bootstrap with \emph{shared} draws across
policies on the same resampled families, so paired contrasts use the same
resample.  Batch top-$q$ decisions couple roots within a batch; where such a
policy is evaluated, resampling uses independent batches or a stated empirical
bootstrap and recomputes the permitted context-only allocation.  A bootstrap that
recomputes allocation is empirical inference, not a replacement safety theorem,
and no per-root i.i.d.\ certificate is claimed.  Rare high-value events, small
$q$ and concentrated support can make effective sample sizes much smaller than
root counts; selected roots and distinct template families per budget and
stratum are reported.
